\documentclass[10pt,landscape]{article}
\usepackage[margin=1.6cm]{geometry}
\usepackage[utf8]{inputenc}
\usepackage[T1]{fontenc}
\usepackage{amsmath}
\usepackage[spanish,es-noshorthands]{babel}
\usepackage{booktabs}
\usepackage{float}
\usepackage{needspace}
\usepackage{tabularx}
\usepackage{longtable}
\usepackage{array}
\usepackage{xcolor}
\usepackage{colortbl}
\usepackage{ragged2e}
\definecolor{prim}{HTML}{1B2A4A}
\definecolor{primlight}{HTML}{D6E4F0}
\definecolor{gris}{HTML}{F7F7F5}
\newcolumntype{L}[1]{>{\RaggedRight\arraybackslash}p{#1}}
\renewcommand{\arraystretch}{1.12}
\setlength{\aboverulesep}{0pt}
\setlength{\belowrulesep}{0pt}
\title{\textcolor{prim}{Planificación LAB019--021: réplica de LAB004 · LAB006 · LAB009}\\[2pt]\large Vendimia 2026 --- mosto natural, escala laboratorio (v2)}
\author{}
\date{22/09/2026 (v2: muestreo 09:00/12:00/15:00 · perfiles en múltiplos de 12 h · F3 = Protocolo B real)}
\begin{document}
\maketitle
\thispagestyle{empty}
\vspace{-1.2cm}

\section{Objeto y contexto}
Se replican tres fermentaciones históricas sobre el mismo diseño (mosto natural \emph{Sauvignon blanc}, levadura Zymaflore X5, reactores R-2L de 2 L): \textbf{LAB004} (Protocolo A, 06/04) en F1, \textbf{LAB006} (Protocolo C, 06/04) en F2 y \textbf{LAB009} (Protocolo B real, 17/04) en F3, cuya medición de CO$_2$ presentó fugas/perfil no confiable y fue descartada del análisis (\texttt{EXCLUDED\_BATCHES}). Con el sistema de CO$_2$ correctamente conectado y verificado, se repiten los perfiles de temperatura y nutrición para obtener una serie continua confiable, y se agrega la medición de CO$_2$ disuelto con \textbf{Carbodoseur} en eventos clave (observabilidad del pool de CO2 disuelto del modelo).

\textbf{Correspondencia:} LAB019 = R-2L-01 (F1), Protocolo A, réplica de LAB004; LAB020 = R-2L-02 (F2), Protocolo C, réplica de LAB006; LAB021 = R-2L-03 (F3), Protocolo B, réplica de LAB009 (grupo del 17/04, sin registro en Maestro). Los tres experimentos originales tuvieron CO$_2$ descartado por QC, por eso son los que más ganan con esta réplica.

\textbf{Decisión de diseño (21/09):} setpoints y pulsos 2 en \textbf{múltiplos de 12 h} desde t$_0$ (16:00) $\Rightarrow$ \emph{todos} los eventos automáticos caen a las 16:00 (con personal presente, salvo el sábado). Redondeos respecto del original: F1 99$\rightarrow$96 y 171$\rightarrow$168; F2 123$\rightarrow$120 y 219$\rightarrow$216; pulsos 51$\rightarrow$48 y 75$\rightarrow$72; F3 $\sim$68$\rightarrow$72.

\textbf{Nota de auditoría (LAB009):} la homologación (\texttt{Perfiles\_temperatura}) registra LAB009 como 18$\rightarrow$21\,\textdegree C a t=96 h$\rightarrow$16\,\textdegree C a t=168 h, pero su \emph{logger crudo} muestra inicio en 21\,\textdegree C y caída a 18\,\textdegree C a t$\approx$68 h; las etiquetas A/B de LAB007/LAB009 aparecen intercambiadas en \texttt{00\_Resumen} y el cambio final de LAB006 (t=219 h) tampoco quedó en el logger. $\theta_{\text{natural}}$ calibró con la temperatura invertida de LAB009 (corregir el dataset antes de re-calibrar). Esta réplica programa el perfil real (B).

\section{Análisis de los experimentos originales}
De las hojas \texttt{Eventos\_operacion}, \texttt{Perfiles\_temperatura} y \texttt{LAB004..LAB006, LAB009} del workbook homologado + loggers crudos (\texttt{raw\_data}):

\begin{table}[h]\centering
\begin{tabular}{L{3.6cm} L{4.4cm} L{4.6cm} L{4.6cm}}
\toprule
\rowcolor{primlight}\textbf{Hito} & \textbf{A --- LAB019 (de LAB004)} & \textbf{B --- LAB021 (de LAB009)} & \textbf{C --- LAB020 (de LAB006)} \\
\midrule
Setpoint inicial & 18\,\textdegree C & 21\,\textdegree C (logger; homologación decía 18) & 16\,\textdegree C \\
Pulso nutricional 2 & t=48 h (orig.: 51 h, densidad 1040) & t=48 h (ICS; sin registro) & t=72 h (orig.: 75 h, densidad 1040) \\
Cambio de T 1 & 21\,\textdegree C a t=96 h (orig.: 99) & 18\,\textdegree C a t=72 h (logger orig.: $\approx$68) & 18\,\textdegree C a t=120 h (orig.: 123) \\
Cambio de T 2 & 16\,\textdegree C a t=168 h (orig.: 171) & \emph{omitido} (sin evidencia de ejecución) & 21\,\textdegree C a t=216 h (orig.: 219) \\
Fermentación seca (GF$<$4 g/L) & t$\approx$122--146 h & t$\approx$120--144 h & t$\approx$170--194 h \\
Última muestra del original & t=163 h & t=160 h & t=211 h (15 muestras) \\
Duración total & $\approx$163 h ($\approx$6.8 días) & $\approx$160 h ($\approx$6.7 días) & $\approx$211 h ($\approx$8.8 días) \\
\bottomrule
\end{tabular}
\caption{Hitos de los experimentos originales (LAB004 y LAB006: siembra 06/04/2026 13:00; LAB009: siembra 17/04/2026 16:00) y su traducción a múltiplos de 12 h.}
\end{table}

Con $t_0$ = lunes 21/09 a las 16:00: \textbf{A} seca entre el sábado 26 16:00 y el domingo 27 18:00 (t=120--146 h), último cambio T de F1 el lunes 28/09 a las 16:00 (t=168 h) y F3 en 18\,\textdegree C desde el jueves 24 (t=72 h); \textbf{C} seca entre el lunes 28 y el martes 29/09 por la tarde (t=170--194 h), último cambio el miércoles 30/09 a las 16:00 (t=216 h). \textbf{La campaña completa cierra el jueves 01/10/2026} (muestra condicional de F2 + desmontaje).

\section{Anclas temporales y definiciones}
\begin{itemize}\setlength{\itemsep}{1pt}
\item \textbf{PI (pre-inóculo):} muestra del mosto antes de inocular, lunes 15:30. Sin Oculyze. Cierra la deuda de $E_0$/YAN$_0$ y da la condición inicial del pool de CO$_2$ disuelto (en agosto el mosto arrancó con $\approx$430--470 mg/L, mientras el modelo hoy inicializa el pool en 0).
\item \textbf{t$_0$:} lunes 21/09/2026 \textbf{16:00}, fin del pulso de inoculación de la bomba peristáltica (programada 15:55$\rightarrow$16:00, $\approx$4--5 min). Todos los tiempos $t$ [h] se miden desde t$_0$.
\item \textbf{Muestreo:} mar--jue \textbf{09:00, 12:00 y 15:00} (3 muestras/día) + \textbf{intermedia opcional} el martes 16:45; viernes \textbf{09:00 + 11:00}. El par PI/t$_0$ (15:30 / 16:10) queda a 40 min: única excepción corta junto con la intermedia (1.75 h, mini).
\item \textbf{Visitas y procesamiento:} visita completa (09:00 y 15:00) demora \textbf{$\approx$1.5--2 h} (3 fermentadores $\times$ Y15, Oculyze, Carbodoseur, Brix, densidad, DO); la de \textbf{12:00 es mini} ($\approx$15--20 min, solo instrumental) para no exceder la densidad analítica de los originales (2 paneles/día) y proteger el almuerzo. Viernes: salida 12:30. Fin de semana sin personal.
\end{itemize}

\Needspace*{20\baselineskip}
\section{Perfiles de temperatura (programados en baños/controladores)}
\begin{table}[H]\centering
\begin{tabular}{L{3.4cm} L{1.8cm} L{2.6cm} L{1.8cm} L{3.1cm} L{7.6cm}}
\toprule
\rowcolor{primlight}\textbf{Reactor} & \textbf{Fase} & \textbf{Setpoint} & \textbf{t [h]} & \textbf{Fecha-hora} & \textbf{Nota} \\
\midrule
LAB019 (A, de LAB004) & 1 & 18 \textdegree C & 0 & lun 21/09 16:00 & Setpoint inicial; mosto cargado ya equilibrado \\LAB019 (A, de LAB004) & 2 & 21 \textdegree C & 96 & vie 25/09 16:00 & Automático, con personal (original: 99 h) \\LAB019 (A, de LAB004) & 3 & 16 \textdegree C & 168 & lun 28/09 16:00 & Automático, con personal (original: 171 h) \\LAB020 (C, de LAB006) & 1 & 16 \textdegree C & 0 & lun 21/09 16:00 & Setpoint inicial \\LAB020 (C, de LAB006) & 2 & 18 \textdegree C & 120 & s\'ab 26/09 16:00 & Automático, fin de semana (original: 123 h) \\LAB020 (C, de LAB006) & 3 & 21 \textdegree C & 216 & mi\'e 30/09 16:00 & Automático, con personal (original: 219 h) \\LAB021 (B, de LAB009) & 1 & 21 \textdegree C & 0 & lun 21/09 16:00 & Setpoint inicial (perfil real por logger; la homologación decía 18 \textdegree C) \\LAB021 (B, de LAB009) & 2 & 18 \textdegree C & 72 & jue 24/09 16:00 & Automático, con personal (logger original: caída a $\approx$68 h); mismo instante que el pulso de F2 \\
\bottomrule
\end{tabular}
\caption{Tiempos en múltiplos de 12 h desde t$_0$ (decisión 21/09), redondeados desde los ejecutados en abril. \emph{F3 (B nominal): el segundo cambio 21$\rightarrow$16\,\textdegree C se omite --- sin evidencia de que LAB009 lo ejecutara (desviación documentada).}}
\end{table}

\section{Nutrición (automatizada con bomba peristáltica)}
\begin{itemize}\setlength{\itemsep}{1pt}
\item \textbf{Pulso 1 (con la inoculación):} Springferm Xtrem 1.00 g + FDA 0.40 g por reactor, \emph{dentro del inóculo} (50\% del suplemento total; igual en A, B y C). Lunes 15:55$\rightarrow$16:00.
\item \textbf{Pulso 2 F1+F3 (simultáneo):} igual dosis, t=48 h $\Rightarrow$ \textbf{mié 23/09 16:00} (originales: LAB004 51 h con densidad 1040; LAB009 48 h vía ICS). \emph{Contingencia (visita -08, mié 15:00):} F1 esperada $\sim$1044--1046 g/L (LAB004: 1044.3 a t=50 h) y F3 $\sim$1039 g/L; si F1 $\geq$1052 $\rightarrow$ reprogramar a jue 09:00; si $\leq$1036 $\rightarrow$ ejecutar en la misma visita.
\item \textbf{Pulso 2 F2 (C, de LAB006):} t=72 h $\Rightarrow$ \textbf{jue 24/09 16:00} (original: 75 h; LAB006: densidad 1040.0 g/L a t=74 h); mismo instante que el cambio T de F3. \emph{Contingencia (visita -11, jue 15:00):} esperada $\sim$1040--1042 g/L; si $\geq$1048 $\rightarrow$ vie 09:00; si $\leq$1032 $\rightarrow$ ejecutar en la misma visita.
\item Verificar en el log de la bomba (flanco de caída) la ejecución de cada pulso a la mañana siguiente (jue: F1/F3; vie: F2).
\end{itemize}

\section{Plan de muestreo}
Mar--jue: \textbf{3 muestras/día} a las \textbf{09:00} (visita completa), \textbf{12:00} (mini instrumental: Carbodoseur + densidad + Brix + DO, sin Y15/Oculyze/Alcolyzer, $\approx$15--20 min) y \textbf{15:00} (visita completa), más una \textbf{intermedia opcional} el martes 16:45 (mini; punto de onset). Viernes: \textbf{09:00} completa + \textbf{11:00} mini (1--2 muestras según sequedad; salida 12:30). Separación 3 h intra-día (16--18 h entre días). Análisis por visita completa: \textbf{Y15} (glucosa, fructosa, PAN, NH$_4\rightarrow$YAN, glicerol), \textbf{Alcolyzer} (etanol), \textbf{Oculyze} (queda en su software), \textbf{densidad + Brix + DO}. Semana siguiente: cola liviana (lun 28 $\times$2, mar 29 $\times$1, mié 30 $\times$2, jue 01 condicional).

\subsection*{Carbodoseur (CO$_2$ disuelto) — 10 lecturas fijas + 1 condicional}
\begin{itemize}\setlength{\itemsep}{1pt}
\item \textbf{Eventos:} PI, -01; martes -02, -03, -04 y -05 opc (17 / 20 / 23 / 24.75 h: trayectoria y onset); miércoles -06 y -08 (41 y 47 h: onset de C y meseta); jueves -09 (65 h: post-pulso F1/F3); viernes -12 (89 h: post-pulso F2); lunes 28 -14 (161 h: contrastes Csat(T)); -19 cond (233 h: post-salto 21\,\textdegree C de F2) ($\approx$100 mL por lectura; la alícuota se devuelve al reactor).
\item \textbf{Por qué estos puntos:} IC del pool en PI/t0 (el modelo hoy parte en 0 y el mosto real arrancó con $\sim$450 mg/L); \textbf{cuatro puntos el martes} cubren la subida del pool hasta el onset (B $\sim$22--24 h ref.\ LAB016--018; A 27--33 h ref.\ loggers de abril), la ventana donde \emph{threshold} vs \emph{continuous release} más divergen; miés: onset de la C fría y meseta vs Csat; post-pulsos (jue/vie) para la respuesta al pulso; -14 tras los cambios T para validar Csat(T).
\item \textbf{Protocolo:} tomar SIEMPRE la primera alícuota, directo del puerto y sin burbujear; anotar \emph{volumen remanente} (mL) y T de medición en la planilla de registro; el CO$_2$ disuelto se completa después con la tabla digitalizada. Leer DO en la misma visita. PI sale del contenedor de mosto (no consume volumen de reactores).
\item \textbf{Volumen (balance real):} la alícuota del Carbodoseur ($\approx$100 mL) \textbf{se devuelve al reactor}: no se pierde. Las únicas pérdidas no retornables por reactor son: Alcolyzer \textbf{50 mL por muestra} (F1/F3: mar--mié--jue $\approx$3 muestras = 150 mL; F2 = LAB020, la C fría: $\approx$5 muestras = 250 mL); Y15 1 mL y Oculyze 1 mL por visita completa; $\approx$2 mL por visita de mangueras/operación (también en minis). Pérdida total $\approx$0.2 L (F1/F3) y $\approx$0.3 L (F2): volumen final $\approx$1.7--1.8 L de los 2.0 L iniciales. Registrar el volumen acumulado solo como control.
\end{itemize}

{\footnotesize
\begin{longtable}{c L{2.2cm} L{1.6cm} c L{1.9cm} L{1.9cm} L{1.9cm} c c L{5.8cm}}
\toprule
\rowcolor{primlight}\textbf{N$^\circ$} & \textbf{Fecha-hora} & \textbf{Día} & \textbf{t [h]} & \textbf{LAB019 (F1)} & \textbf{LAB020 (F2)} & \textbf{LAB021 (F3)} & \textbf{Carbo} & \textbf{Sep.} & \textbf{Nota} \\
\midrule
\endfirsthead
\toprule
\rowcolor{primlight}\textbf{N$^\circ$} & \textbf{Fecha-hora} & \textbf{Día} & \textbf{t [h]} & \textbf{LAB019 (F1)} & \textbf{LAB020 (F2)} & \textbf{LAB021 (F3)} & \textbf{Carbo} & \textbf{Sep.} & \textbf{Nota} \\
\midrule
\endhead
1 & 21/09 15:30 & lun & -- & PI & PI & PI & SÍ & -- & Pre-inóculo: cierra la deuda de E0/YAN0 y da la IC del pool de CO2 disuelto ($\neq$0). No consume volumen de los reactores \\2 & 21/09 16:10 & lun & 0.17 & 01 & 01 & 01 & SÍ & 0.67* & Muestra t=0. Par ancla PI/t0 (40 min): excepción documentada a la separación mínima; YAN de t0 resuelve si el pulso 1 entra o no en la química; X0 por Oculyze \\3 & 22/09 09:00 & mar & 17 & 02 & 02 & 02 & SÍ & 16.83 & Primera del martes (t=17 h): pool pre-onset. Revisar trazas CO2/T de la noche (sin fugas ni saltos) \\4 & 22/09 12:00 & mar & 20 & 03 & 03 & 03 & SÍ & 3 & Mini-muestreo (t=20 h): trayectoria del pool pre-onset; densidad + Brix + DO + Carbo (\textasciitilde{}15-20 min) \\5 & 22/09 15:00 & mar & 23 & 04 & 04 & 04 & SÍ & 3 & Frontera del onset (t=23 h): F3 (B, 21 \textdegree{}C) debería estar subiendo (ref. 22-24 h en LAB016-018); F1 (A) aún pre-onset (ref. abril: 27-33 h) \\6 & 22/09 16:45 & mar & 24.75 & 05 & 05 & 05 & SÍ & 1.75 & INTERMEDIA OPCIONAL (t=24.75 h): en/frontera del onset de F1/F3; el punto que discrimina acumulación previa (threshold) vs emisión temprana (continuous) \\7 & 23/09 09:00 & mié & 41 & 06 & 06 & 06 & SÍ & 16.25 & Post-noche (t=41 h): onset de C (F2, arranque frío) + meseta de A/B; contrastar pool con trazas nocturnas \\8 & 23/09 12:00 & mié & 44 & 07 & 07 & 07 & -- & 3 & Mini-muestreo de densidad (vigilar cruce de 1040 para el pulso de F1/F3 de mañana 16:00) \\9 & 23/09 15:00 & mié & 47 & 08 & 08 & 08 & SÍ & 3 & Pre-pulso 2 de F1/F3 (16:00, t=48 h). Contingencia densidad: F1 esperada \textasciitilde{}1044-1046 g/L (LAB004: 1044.3 a t=50); F3 \textasciitilde{}1039 (LAB009 cruzó 1040 \textasciitilde{}t=44). Si F1 $\geq$1052 reprogramar a jue 09:00; si $\leq$1036 ejecutar en esta visita \\10 & 24/09 09:00 & jue & 65 & 09 & 09 & 09 & SÍ & 18 & Post-pulso F1/F3 (17 h). Verificar logs de bomba del miércoles. F2 pre-pulso (72 h = hoy 16:00) \\11 & 24/09 12:00 & jue & 68 & 10 & 10 & 10 & -- & 3 & Mini-muestreo de densidad (vigilar cruce de 1040 de F2 para el pulso de las 16:00) \\12 & 24/09 15:00 & jue & 71 & 11 & 11 & 11 & -- & 3 & Pre-pulso de F2 (16:00, t=72 h). Contingencia: esperada \textasciitilde{}1040-1042 g/L (LAB006: 1040.0 a t=74); si $\geq$1048 reprogramar a vie 09:00; si $\leq$1032 ejecutar en esta visita. F3 cambia a 18 \textdegree{}C hoy 16:00 (con personal) \\13 & 25/09 09:00 & vie & 89 & 12 & 12 & 12 & SÍ & 18 & Post-pulso F2 (17 h) + post-cambio T de F3 (17 h a 18 \textdegree{}C). Verificar logs de ayer (pulso F2 + setpoint F3). F1 cambia a 21 \textdegree{}C hoy 16:00 \\14 & 25/09 11:00 & vie & 91 & 13 & 13 & 13 & -- & 2 & Última del día antes del cambio T de F1 (16:00); cierre de semana (salida 12:30) \\15 & 28/09 09:00 & lun & 161 & 14 & 14 & 14 & SÍ & 70 & Post-fin de semana (t=161 h): seca esperada de A (t=120-146 h cayó en la ventana ciega; el CO2 continuo la cubre). F1 aún a 21 \textdegree{}C, F2/F3 a 18 \textdegree{}C; F1 cambia a 16 \textdegree{}C hoy 16:00. Descargar/respaldar datos \\16 & 28/09 15:00 & lun & 167 & 15 & 15 & 15 & -- & 6 & Confirmación de sequedad de A: F1/F3 secas si GF<2 g/L en -14/-15 (entonces vaciar el martes) \\17 & 29/09 09:00 & mar & 185 & 16 & 16 & 16 & -- & 18 & Confirmación A (si -15 no bastó) + cola de C. Si A confirmada: vaciar/desmontar F1 y F3 hoy \\18 & 30/09 09:00 & mié & 209 & 17 & 17 & 17 & -- & 24 & C en fase final (LAB006: GF=4.7 g/L a t=170 h). F2 cambia a 21 \textdegree{}C hoy 16:00 (con personal) \\19 & 30/09 15:00 & mié & 215 & 18 & 18 & 18 & -- & 6 & Pre-cambio T de F2 (16:00): C casi seca; base para decidir la muestra condicional de mañana \\20 & 01/10 09:00 & jue & 233 & --- & -19 & --- & cond & 18 & CONDICIONAL: solo F2 y solo si GF(-18) > 2 g/L; 17 h después del salto a 21 \textdegree{}C (cierre del par térmico Csat(T)). Luego vaciado de F2 y cierre de campaña \\
\bottomrule
\end{longtable}
}
\footnotesize *Par ancla PI/t$_0$ (excepción documentada). Sep.\ = horas desde la visita anterior; 1.75 h = intermedia opcional (mini). Carbo = lectura de Carbodoseur (CO$_2$ disuelto); cond = solo si se toma -19.

\subsection*{Criterios de término}
\begin{itemize}\setlength{\itemsep}{1pt}
\item \textbf{A (F1/F3 = LAB019/021):} GF $<$ 2 g/L en dos muestras consecutivas (esperado -14/-15, lunes 28/09); el martes 29 por la mañana se vacían y desmontan F1/F3 si se confirma.
\item \textbf{C (F2 = LAB020):} GF $<$ 2 g/L en dos consecutivas (esperado -17/-18, miércoles 30/09); -19 (jueves 01/10 09:00) solo si GF(-18) $>$ 2 g/L; luego vaciado de F2 y cierre.
\item \textbf{Ventana ciega de fin de semana:} la sequedad de A (t=120--146 h) caerá entre el sábado 16:00 y el domingo 18:00 sin muestreo; la señal continua de CO$_2$ ---justo lo que esta réplica corrige--- cubre el intervalo.
\end{itemize}

\section{Cronograma día a día}
{\footnotesize
\renewcommand{\arraystretch}{1.15}
\begin{longtable}{L{1.7cm} L{1.2cm} L{1.1cm} c L{4.6cm} L{1.6cm} L{1.7cm} L{7.2cm}}
\toprule
\rowcolor{primlight}\textbf{Fecha} & \textbf{Día} & \textbf{Hora} & \textbf{t [h]} & \textbf{Evento} & \textbf{Reactores} & \textbf{Tipo} & \textbf{Detalle} \\
\midrule
\endfirsthead
\toprule
\rowcolor{primlight}\textbf{Fecha} & \textbf{Día} & \textbf{Hora} & \textbf{t [h]} & \textbf{Evento} & \textbf{Reactores} & \textbf{Tipo} & \textbf{Detalle} \\
\midrule
\endhead
21/09 & lun & 08:30 &  & Verificación del sistema CO2 sin fugas & F1,F2,F3 & Preparación & Prueba de presión/agua jabonosa en todas las conexiones y trampas; zero-check de sensores; comparar baseline F1 vs F2/F3 \\21/09 & lun & 10:00 &  & Calibración de bombas peristálticas & F1,F2,F3 & Preparación & Verificar caudal/volumen de las bombas de inóculo y nutrición; cargar programas de la semana (pulsos 48/48/72 h) y setpoints de baños (96/120/72/168/216 h) \\21/09 & lun & 10:00 &  & Baños/controladores a setpoint inicial & F1,F2,F3 & Preparación & F1: 18 \textdegree{}C; F2: 16 \textdegree{}C; F3: 21 \textdegree{}C (para que el mosto llegue equilibrado a t0) \\21/09 & lun & 11:00 &  & Preparación de dosis, inóculo y material & F1,F2,F3 & Preparación & Pesar Springferm Xtrem (1.00 g/dosis) y FDA (0.40 g/dosis); rehidratar inóculo Zymaflore X5; tubos, etiquetas, hielo; Y15 y Alcolyzer operativos; imprimir planilla de registro; iniciar loggers CO2/T \\21/09 & lun & 14:30 &  & Carga y homogeneización del mosto & F1,F2,F3 & Preparación & Mosto natural Sauvignon blanc, 2 L por reactor; registrar volumen y temperatura \\21/09 & lun & 15:30 &  & MUESTRA PI (pre-inóculo) + Carbodoseur & F1,F2,F3 & Muestreo & Carbodoseur primero (\textasciitilde{}100 mL/reactor), luego panel. Y15 de PI puede procesarse mientras corre la inoculación \\21/09 & lun & 15:40 &  & Carga de inóculo + suplemento en línea de bomba & F1,F2,F3 & Preparación & El suplemento del pulso 1 va dentro del inóculo \\21/09 & lun & 15:55 &  & Inoculación automatizada (pulso 1) & F1,F2,F3 & Automático & Bomba peristáltica 15:55$\rightarrow$16:00 (\textasciitilde{}4-5 min, ref. LAB016-018: 232-246 s) \\21/09 & lun & 16:00 & 0 & t0 = fin de la inoculación & F1,F2,F3 & Automático & Ancla temporal de toda la campaña \\21/09 & lun & 16:10 & 0.17 & MUESTRA -01 (t=0) + Carbodoseur & F1,F2,F3 & Muestreo & Panel completo; lunes ajustado (PI y t0 en \textasciitilde{}2 h): si no alcanza, Y15 restante en la tarde-jueves \\21/09 & lun & 16:15 & 0.25 & Verificación inicial de trazas & F1,F2,F3 & Verificación & Confirmar que CO2/T están registrando y que la bomba logueó la inoculación; 17:00-17:30 limpieza diaria \\22/09 & mar & 09:00 & 17 & MUESTRA 02 02 (t=17 h) + Carbodoseur & F1,F2,F3 & Muestreo & Primera del martes (t=17 h): pool pre-onset. Revisar trazas CO2/T de la noche (sin fugas ni saltos) \\22/09 & mar & 12:00 & 20 & MINI-MUESTRA 03 (t=20 h) + Carbodoseur & F1,F2,F3 & Muestreo & Mini-muestreo (t=20 h): trayectoria del pool pre-onset; densidad + Brix + DO + Carbo (\textasciitilde{}15-20 min) \\22/09 & mar & 15:00 & 23 & MUESTRA 04 04 (t=23 h) + Carbodoseur & F1,F2,F3 & Muestreo & Frontera del onset (t=23 h): F3 (B, 21 \textdegree{}C) debería estar subiendo (ref. 22-24 h en LAB016-018); F1 (A) aún pre-onset (ref. abril: 27-33 h) \\22/09 & mar & 16:45 & 24.75 & MINI-MUESTRA (opcional) 05 (t=24.75 h) + Carbodoseur & F1,F2,F3 & Muestreo & INTERMEDIA OPCIONAL (t=24.75 h): en/frontera del onset de F1/F3; el punto que discrimina acumulación previa (threshold) vs emisión temprana (continuous) \\23/09 & mié & 09:00 & 41 & MUESTRA 06 06 (t=41 h) + Carbodoseur & F1,F2,F3 & Muestreo & Post-noche (t=41 h): onset de C (F2, arranque frío) + meseta de A/B; contrastar pool con trazas nocturnas \\23/09 & mié & 12:00 & 44 & MINI-MUESTRA 07 (t=44 h) & F1,F2,F3 & Muestreo & Mini-muestreo de densidad (vigilar cruce de 1040 para el pulso de F1/F3 de mañana 16:00) \\23/09 & mié & 15:00 & 47 & MUESTRA 08 08 (t=47 h) + Carbodoseur & F1,F2,F3 & Muestreo & Pre-pulso 2 de F1/F3 (16:00, t=48 h). Contingencia densidad: F1 esperada \textasciitilde{}1044-1046 g/L (LAB004: 1044.3 a t=50); F3 \textasciitilde{}1039 (LAB009 cruzó 1040 \textasciitilde{}t=44). Si F1 $\geq$1052 reprogramar a jue 09:00; si $\leq$1036 ejecutar en esta visita \\23/09 & mié & 16:00 & 48 & Pulso nutricional 2 — F1 y F3 (simultáneo) & F1,F3 & Automático & Springferm Xtrem 1.00 g + FDA 0.40 g c/u; t=48 h (redondeo de 51/48 h del original). Llegando de la visita de las 15:00 \\24/09 & jue & 09:00 & 65 & MUESTRA 09 09 (t=65 h) + Carbodoseur & F1,F2,F3 & Muestreo & Post-pulso F1/F3 (17 h). Verificar logs de bomba del miércoles. F2 pre-pulso (72 h = hoy 16:00) \\24/09 & jue & 09:00 & 65 & Verificar log: pulsos 2 de F1 y F3 ejecutados & F1,F3 & Verificación & Revisar flancos de la bomba en el logger (análogo a nutricion\_activa de LAB016-018) \\24/09 & jue & 12:00 & 68 & MINI-MUESTRA 10 (t=68 h) & F1,F2,F3 & Muestreo & Mini-muestreo de densidad (vigilar cruce de 1040 de F2 para el pulso de las 16:00) \\24/09 & jue & 15:00 & 71 & MUESTRA 11 11 (t=71 h) & F1,F2,F3 & Muestreo & Pre-pulso de F2 (16:00, t=72 h). Contingencia: esperada \textasciitilde{}1040-1042 g/L (LAB006: 1040.0 a t=74); si $\geq$1048 reprogramar a vie 09:00; si $\leq$1032 ejecutar en esta visita. F3 cambia a 18 \textdegree{}C hoy 16:00 (con personal) \\24/09 & jue & 16:00 & 72 & Pulso nutricional 2 — F2 & F2 & Automático & Springferm Xtrem 1.00 g + FDA 0.40 g; t=72 h (redondeo del original 75 h; LAB006 cruzó 1040 g/L a t=74 h) \\24/09 & jue & 16:00 & 72 & Setpoint 18 \textdegree{}C — F3 & F3 & Automático & 21->18 \textdegree{}C a t=72 h (logger del original: caída a t$\approx$68 h); mismo instante que el pulso de F2, con personal \\25/09 & vie & 09:00 & 89 & MUESTRA 12 12 (t=89 h) + Carbodoseur & F1,F2,F3 & Muestreo & Post-pulso F2 (17 h) + post-cambio T de F3 (17 h a 18 \textdegree{}C). Verificar logs de ayer (pulso F2 + setpoint F3). F1 cambia a 21 \textdegree{}C hoy 16:00 \\25/09 & vie & 09:00 & 89 & Verificar log: pulso de F2 y setpoint de F3 (jue 16:00) & F2,F3 & Verificación & Ídem \\25/09 & vie & 11:00 & 91 & MINI-MUESTRA 13 (t=91 h) & F1,F2,F3 & Muestreo & Última del día antes del cambio T de F1 (16:00); cierre de semana (salida 12:30) \\25/09 & vie & 12:00 & 92 & Cierre de fin de semana & F1,F2,F3 & Verificación & Confirmar setpoint F1 21 \textdegree{}C ejecutado (16:00) y F2 18 \textdegree{}C programado (sáb 16:00); capacidad de loggers hasta el lunes; 12:00-12:30 limpieza y salida \\25/09 & vie & 16:00 & 96 & Setpoint 21 \textdegree{}C — F1 & F1 & Automático & 18->21 \textdegree{}C a t=96 h (redondeo del original 99 h); viernes 16:00, con personal \\\rowcolor{gris} 26/09 & sáb & 00:00 &  & Fin de semana sin personal & F1,F2,F3 & Sin personal & Sábado y domingo: logging continuo CO2/T; único evento: setpoint F2 18 \textdegree{}C (sáb 16:00, automático). Ventana ciega de muestreo: la sequedad de A (t=120-146 h) ocurrirá aquí; el CO2 continuo sin fugas cubre el intervalo \\\rowcolor{gris} 26/09 & sáb & 16:00 & 120 & Setpoint 18 \textdegree{}C — F2 & F2 & Automático & 16->18 \textdegree{}C a t=120 h (redondeo del original 123 h); sábado, sin personal \\28/09 & lun & 09:00 & 161 & MUESTRA 14 14 (t=161 h) + Carbodoseur & F1,F2,F3 & Muestreo & Post-fin de semana (t=161 h): seca esperada de A (t=120-146 h cayó en la ventana ciega; el CO2 continuo la cubre). F1 aún a 21 \textdegree{}C, F2/F3 a 18 \textdegree{}C; F1 cambia a 16 \textdegree{}C hoy 16:00. Descargar/respaldar datos \\28/09 & lun & 09:00 & 161 & Verificar cambios T del finde + respaldo & F1,F2,F3 & Verificación & Confirmar setpoint 18 \textdegree{}C en F2 (sáb 16:00); F1 cambia a 16 \textdegree{}C hoy 16:00 (con personal); descargar y respaldar datos \\28/09 & lun & 15:00 & 167 & MUESTRA 15 15 (t=167 h) & F1,F2,F3 & Muestreo & Confirmación de sequedad de A: F1/F3 secas si GF<2 g/L en -14/-15 (entonces vaciar el martes) \\28/09 & lun & 16:00 & 168 & Setpoint 16 \textdegree{}C — F1 & F1 & Automático & 21->16 \textdegree{}C a t=168 h (redondeo del original 171 h); lunes 16:00, con personal \\29/09 & mar & 09:00 & 185 & MUESTRA 16 16 (t=185 h) & F1,F2,F3 & Muestreo & Confirmación A (si -15 no bastó) + cola de C. Si A confirmada: vaciar/desmontar F1 y F3 hoy \\29/09 & mar & 09:00 & 185 & Si A confirmada seca (-14/-15): vaciar F1 y F3 & F1,F3 & Cierre & Criterio: GF<2 g/L en 2 muestras consecutivas (-14/-15); desmontaje de reactores liberados \\30/09 & mié & 09:00 & 209 & MUESTRA 17 17 (t=209 h) & F1,F2,F3 & Muestreo & C en fase final (LAB006: GF=4.7 g/L a t=170 h). F2 cambia a 21 \textdegree{}C hoy 16:00 (con personal) \\30/09 & mié & 15:00 & 215 & MUESTRA 18 18 (t=215 h) & F1,F2,F3 & Muestreo & Pre-cambio T de F2 (16:00): C casi seca; base para decidir la muestra condicional de mañana \\30/09 & mié & 16:00 & 216 & Setpoint 21 \textdegree{}C — F2 & F2 & Automático & 18->21 \textdegree{}C a t=216 h (redondeo del original 219 h); miércoles 16:00, con personal \\01/10 & jue & 09:00 & 233 & MUESTRA -19 (condicional, solo F2) + Carbodoseur & F2 & Muestreo & Solo si GF(-18) > 2 g/L; luego vaciado de F2 \\01/10 & jue & 11:00 & 235 & Cierre de campaña & F1,F2,F3 & Cierre & Vaciado/desmontaje de F2, descarga y respaldo final de datos, limpieza del laboratorio \\
\bottomrule
\end{longtable}
}

\section{Restricciones y contingencias}
\begin{itemize}\setlength{\itemsep}{1pt}
\item \textbf{Horarios y procesamiento:} lun--jue 8:00--17:30: muestreos \textbf{09:00, 12:00 (mini) y 15:00} (completas $\approx$1.5--2 h; mini $\approx$15--20 min) + intermedia opcional mar 16:45; almuerzo 13:00--14:30; limpieza desde 17:00. Viernes 8:00--12:30 (09:00 completa + 11:00 mini). Fin de semana sin personal.
\item \textbf{Separación entre muestras:} mínimo 2 h (piso 1.5 h); el grid mar--jue usa 3 h (la de 12:00 es mini, sin Y15). Excepciones documentadas: par PI/t$_0$ (40 min) e intermedia del martes (1.75 h).
\item \textbf{Mosto retrasado:} si no está listo a las 15:30, correr PI y t$_0$ (p.\,ej.\ mosto 16:00 $\rightarrow$ PI 16:00, t$_0$ 16:30) y desplazar \emph{todos} los eventos automáticos por el mismo desfase, manteniendo los tiempos $t$ [h].
\item \textbf{Fugas CO$_2$:} verificación con presión/jabón el lunes a.\,m.; revisar trazas a diario (caída a cero o asimetría F1 vs F2/F3 son señales de fuga).
\item \textbf{Carbodoseur:} primera alícuota de cada visita marcada, sin burbujear; anotar volumen remanente y T; CO$_2$ después con la tabla digitalizada; $\approx$100 mL por lectura \textbf{que se devuelve al reactor} (el «Vol.\ rem.» es solo control operativo). La pérdida real es pequeña: Alcolyzer 50 mL/muestra + Y15/Oculyze 1 mL c/u por visita completa + $\approx$2 mL/visita de mangueras $\Rightarrow$ $\approx$0.2 L (F1/F3) y $\approx$0.3 L (F2) en total; volumen final $\approx$1.7--1.8 L de 2.0 L.
\item \textbf{Señales de alarma en el lag (lecciones de abril):} trazas en 0.000 exactos = línea sospechosa de bloqueo (patrón F1); baseline creciente o saltos = sospecha de fuga (patrón F2/F3); lo esperable es $\sim$0.5--1 sccm plano por desorción.
\item \textbf{Múltiplos de 12:} setpoints y pulsos 2 redondeados a múltiplos de 12 h desde t$_0$ (16:00): todos los eventos automáticos ocurren a las 16:00, con personal presente (salvo el sábado). Diferencias $-3$ h respecto del original ejecutado (99/171, 123/219, 51/75, $\sim$68) quedan documentadas como parte del diseño.
\item \textbf{Trazabilidad:} nombrar canales LAB019/020/021 antes del lunes 16:00; registrar la hora exacta de cada flanco de bomba; descargar y respaldar datos el lunes 28 a.\,m.\ y al cierre.
\end{itemize}

\vspace{4pt}
{\sloppy
\noindent\begin{minipage}{\linewidth}
\textcolor{gray}{\footnotesize Fuentes: \texttt{mosto\_natural\_xthiol.xlsx} (hojas \texttt{Eventos\_operacion}, \texttt{Perfiles\_temperatura}, \texttt{LAB004..LAB006}, \texttt{LAB009}); loggers crudos \texttt{raw\_data} (verificación SP LAB009); \texttt{Planilla\_Maestra\_ID\_2026\_14\_07\_26.xlsx}; \texttt{docs/natural\_must\_temporal\_audit\_LAB001\_018.md}; \texttt{docs/natural\_must\_model\_calibration\_and\_estimator\_readiness\_2026.md} (CO2sat\_scale en cota inferior; Carbodoseur LAB013--015). Registro de mediciones: \texttt{registro\_muestras\_LAB019\_021\_v2\_2026.xlsx}.}
\end{minipage}\par}
\end{document}
